\section{Coeficientes Sísmicos}

Para poder determinar el coeficiente sísmico a utilizar en el análisis sísmico, lo primero que se debe hacer es determinar máximos y mínimos a respetar. Para ello, en la normativa se establece en el artículo 6.2.3.1.1 de la norma NCh 433 \cite{NCh433:1996} que el coeficiente sísmico no debe ser menor a:

\begin{equation}
    C_{min} = \frac{A_0 \cdot S \cdot I}{6} = \frac{0.4 \cdot 0.9 \cdot 1}{6}= 0.0666
\end{equation}

En cuanto al máximo, se establece en el artículo 6.2.3.1.2 que la ecuación del coeficiente sísmico está ubicada en la tabla 6.4 y depende del factor \textbf{R}. Recordando el valor de R calculado en la sección anterior, se obtiene el siguiente valor de coeficiente sísmico máximo:
\begin{equation}
    C_{max} = 0.35 \cdot S \cdot A_0 \cdot I = 0.35 \cdot 0.9 \cdot 0.4 \cdot 1 = 0.126
\end{equation}

Entonces, el coeficiente sísmico a utilizar en el análisis sísmico debe estar entre los valores de $C_{min}$ y $C_{max}$, es decir, entre 0.0666 y 0.126. Para este proyecto, se utilizará un valor de coeficiente sísmico de 0.1, ya que se encuentra dentro del rango establecido y es un valor comúnmente utilizado en proyectos de este tipo.